%% ============================================================
\section{Background}
\label{sec:background}

\subsection{SEI CERT C Coding Standard}

The SEI CERT C Coding Standard is maintained by the Software
Engineering Institute at Carnegie Mellon University.  It defines rules
and recommendations organized into 17 categories, shown in
\cref{tab:cert-categories}.

\begin{table}[t]
\centering
\footnotesize
\begin{tabular}{@{}llr@{}}
\toprule
\textbf{Prefix} & \textbf{Category} & \textbf{Rules} \\
\midrule
API & API Usage & 9 \\
ARR & Arrays & 9 \\
CON & Concurrency & 23 \\
DCL & Declarations & 32 \\
ENV & Environment & 9 \\
ERR & Error Handling & 11 \\
EXP & Expressions & 31 \\
FIO & File I/O & 35 \\
FLP & Floating Point & 14 \\
INT & Integers & 23 \\
MEM & Memory Mgmt & 17 \\
MSC & Miscellaneous & 31 \\
POS & POSIX & 21 \\
PRE & Preprocessor & 16 \\
SIG & Signals & 7 \\
STR & Strings & 16 \\
WIN & Windows & 7 \\
\midrule
& \textbf{Total} & \textbf{311}\textsuperscript{$\dagger$} \\
\bottomrule
\end{tabular}
\caption{CERT C rule categories implemented in SqC.
\textsuperscript{$\dagger$}Includes both mandatory rules and advisory
recommendations.}
\label{tab:cert-categories}
\end{table}

Each rule addresses a specific class of undefined behavior, security
vulnerability, or quality defect.  For example, EXP34-C (``Do not
dereference null pointers'') maps directly to CWE-476; INT32-C
(``Ensure that operations on signed integers do not result in
overflow'') maps to CWE-190.  The standard provides a rule-to-CWE
mapping that enables systematic benchmarking.

\subsection{NIST Juliet Test Suite}

The Juliet Test Suite v1.3~\cite{juliet} (commit \texttt{f88433e})
is developed by NIST's
Software Assurance Metrics and Tool Evaluation (SAMATE) project.
It contains 64,099 C/C++ test cases across 118 CWE categories,
organized into ``good'' functions (correct code that should not
trigger warnings) and ``bad'' functions (code containing the
specified vulnerability).

Juliet provides a controlled evaluation environment: each test case
isolates a single vulnerability type with known ground truth, enabling
precise measurement of true positive (TP) and false positive (FP) rates
per CWE.

Juliet also has known limitations worth acknowledging.  Test cases use
synthetic patterns that may not reflect real-world coding idioms.  The
structured naming conventions (functions prefixed with \texttt{bad} or
\texttt{good}) can inadvertently advantage or disadvantage
pattern-matching tools.  And Juliet uses explicit C types
(\texttt{int}, \texttt{unsigned int}) rather than \texttt{stdint.h}
typedefs (\texttt{uint32\_t}), so fixes targeting typedef resolution
show impact on real-world code but not on Juliet.

These limitations are why we supplement Juliet with real-world
evaluation.

\subsection{Real-World Benchmarks}

Nine open-source C projects serve as real-world benchmarks:
libcrc (1K LOC), lua (31K LOC), raylib (56K LOC), mosquitto (39K LOC),
curl (186K LOC), sqlite (182K LOC), hostap (590K LOC), pure-ftpd (33K
LOC), and seL4 (87K LOC).  These span embedded libraries,
scripting-language runtimes, game/graphics libraries, network servers,
databases, an FTP daemon, and a formally verified microkernel.  Lua and
raylib were added as a fifth and sixth ground-truth oracle specifically
to exercise C99 idioms (compound literals, designated initializers) that
the original five-project set does not contain (\cref{sec:sloc-scope});
pure-ftpd was added as an eighth oracle to exercise SQL-client-API call
sites (\texttt{mysql\_real\_query}, \texttt{PQexec}) absent from the
other codebases, and seL4 as a ninth to test a freestanding,
formally-verified, hardware-facing codebase with no libc/POSIX surface.
